\section{Three missing quaternary columns: a classical-input completion algorithm}
Let $n$ be even and let $V\in\mu_4^{n\times(n-r)}$ satisfy
$V^*V=nI$, with $1\le r\le3$. The following completion statement
uses the \emph{classical real codimension-at-most-seven Hadamard
completion theorem}. That input is explicitly assumed here, not
reproved or claimed new. References include Hall's four-row result
and Verheiden's extension, recorded in
\url{https://thesis.caltech.edu/2899/}, and the primary discussion in
\emph{Counting results for thin Butson matrices},
\url{https://www.combinatorics.org/ojs/index.php/eljc/article/download/v21i3p12/pdf/}.

\begin{lemma}[A one-column extension gate through realification]
For arbitrary $r\ge1$, a quaternary column orthogonal to all columns
of $V$ exists if and only if the real sign frame
$A=\mathcal R((1+i)V)$ has an orthogonal real sign column.
\end{lemma}
\begin{proof}
Write a real sign column in paired coordinates as
$w=(a_1,b_1,\ldots,a_n,b_n)^T$ and let $u_i=a_i+ib_i$.
The paired realified columns give
$A^Tw=0$ if and only if $((1+i)V)^*u=0$, equivalently $V^*u=0$.
Thus $z=u/(1+i)\in\mu_4^n$ is the desired new column.
Conversely $(1+i)z$ has real and imaginary parts in $\{\pm1\}$,
and realifies to such a $w$.
\end{proof}

\begin{theorem}[Automatic completion in complex codimension at most three]
Every such $V$ with $1\le r\le3$ can be completed to a
$\mathrm{BH}(n,4)$ matrix, retaining all its original columns.
\end{theorem}
\begin{proof}
For even $n\ge4$, $A$ has $2n$ rows, $2n-2r$ sign columns and
$A^TA=2nI$. Its real codimension is $2r\le6$, so classical real
completion supplies a full real Hadamard extension and in particular
one new real sign column orthogonal to $A$. The lemma converts that
column to one quaternary column orthogonal to $V$. Append it and
repeat at complex codimension $r-1$. The order is unchanged.
For $n=2$ the assertion follows directly from the order-two matrix:
one prescribed fourth-root column is completed by changing one sign,
and an empty input uses that matrix itself.
\end{proof}
This avoids an unwarranted assumption that an arbitrary supplied real
completion respects the canonical complex structure. Only one real
column is used at each step; its canonical complex-structure partner
is automatically another sign column. The unused real completing
columns may be discarded, and the classical theorem is applied again.

\begin{theorem}[An explicit finite exact algorithm]
For $r\le3$, Gaussian-rational row reduction and at most
$4^{r-1}+\cdots+4+1\le21$ candidate-column tests construct the
completion. The theorem does not require a real Hadamard completion
as algorithmic input.
\end{theorem}
\begin{proof}
Compute a systematic basis of $\ker V^*$: Gaussian elimination
selects $r$ free coordinate positions, and every kernel vector is
determined by its values there. A fourth-root vector can be globally
multiplied by a fourth root so that its first free coordinate equals
one. Its other free coordinates have $4^{r-1}$ possible assignments.
For each assignment compute the unique kernel vector and check all
its entries are in $\mu_4$. The preceding theorem guarantees a
successful assignment. Append that vector and repeat. Its squared
norm is $n$, and its orthogonality to all existing columns is exact,
so the invariants hold at every stage. All coefficients are Gaussian
rational and every test can be done exactly.
\end{proof}
For fixed $r\le3$ the matrix elimination costs $O(n^3)$ scalar
operations; after a kernel basis is available each candidate costs
$O(nr)$ operations. This is an exact-arithmetic operation count,
not an unrestricted bit-complexity or approximate-input claim.

\begin{corollary}[A precise codimension-four gate]
For a quaternary input with $r=4$, full quaternary completion is
equivalent to the existence of even one new fourth-root column.
The latter is decided by at most $4^3=64$ normalized candidate tests.
If a new column exists, the remaining defect three is automatically
completed by the theorem. No automatic existence of that first
column at defect four is asserted.
\end{corollary}
In particular a maximal proper quaternary frame at even order must
miss at least four columns. The result closes a fixed-size residual
gap once a large partial frame is supplied; it does not construct
that input at unknown orders. The completion theory and its small-
alphabet consequence are not advertised as bibliographically new.
